
\def\PUDcodigo{ECA221}

\def\PUDcodacad{04507.35}

\def\PUDdisciplina{Controladores Lógicos Programáveis}

\def\PUDsemestre{S7}

\def\PUDhoras{80}

%NOVOS ---------------------------------------------------
\def\PUDhorasTeoria{40}
\def\PUDhorasPratica{40}

\def\PUDinterECA{ 
\hyperref[PUD:ECA213]{Microcontroladores (04507.28)}

\hyperref[PUD:ECA223]{Máquinas Elétricas (04507.36)}

\hyperref[PUD:ECA228]{A. Hidráulicos e Pneumáticos (04507.37)}
}

\def\PUDatuaisECA{
- Introdução à Industria 4.0

- Aplicações de soluções da Industria 4.0
}

\def\PUDrecursos{Projetor multimídia; Quadro branco e pincel; Aulas em laboratório.}

%----------------------------------------------------------

\def\PUDcreditos{4}

\def\PUDprerequisito{ \hyperref[PUD:ECA211]{ECA211} }

\def\PUDpreacad{ \hyperref[PUD:ECA213]{Microcontroladores (04507.28)} }

\def\PUDobjetivos{Compreender e desenvolver programas para CLP;  Diagnosticar e corrigir falhas em sistemas de automação; Projetar um sistema de controle com uso de CLP. Além de verificar o contexto de aplicações regionais e interdiciplinariedade tanto com as disciplinas de pré-requisito quanto com as disicplinas do semestre corrente, como Máquinas Elétricas, Acionamentos Hidráulicos e Pneumáticos e outras.} %A ÚLTIMA FRASE É NOVA

\def\PUDementa{Conceitos de Automação industrial;  Controladores Lógicos Programáveis (CLP);  Estrutura básica e Princípio de funcionamento de um CLP;  Linguagem de programação conforme norma IEC 61131-3;  Programação de controladores programáveis;  Programação em Ladder;  Noções de sistema SCADA com uso do CLP;  projeto de um sistema de controle com uso do CLP.}

\def\PUDconteudo{ & Controladores Lógicos Programáveis;   Introdução.  Breve histórico.  Evolução.  Aplicações.  Arquiteturas: compacto, modular, I/O distribuído; 
\\ 
 & Arquitetura de CLPs: Micromprocessador, tipos de processamentos (cíclico, interrupção, comandado por tempo e por evento.  Mapa de memória.  Dispositivos de entrada e saída, terminal de programação. 
\\ 
 &  Princípio de funcionamento de um CLP: Estados de operação.  Funcionamento interno do CLP.  Linguagem de programação (baixo nível e alto nível); 
\\ 
 &  Programação de controladores programáveis: Tipos das linguagens de programação ( Ladder diagram (ld) - diagrama de contatos, Function blocks diagram (fbd) - diagrama de blocos, Instruction list (il) - lista de instrução, Structured text (st) – texto estruturado, Sequential function chart (sfc) - passos ou step e Linguagem corrente ou natural).  Normalização - IEC 61131
\\ 
 &  Programação em Ladder: Desenvolvimento do programa Ladder.  Associação de contatos no Ladder.  Instruções básicas. 
\\ 
 &  Noções de sistema SCADA com uso do CLP: Arquitetura da rede clp para sistemas SCADA. Aplicações de soluções da Industria 4.0.
\\ 
 & Desenvolvimento de Projetos de Automação com a utilizaçaõ do CLP. }

\def\PUDmetodo{Aulas expositórias que podem ser teóricas e/ou práticas, onde as práticas no laboratório serão marcadas no decorrer da disciplina.}

\def\PUDavalia{A avaliação será feita com aplicação de provas teóricas e/ou práticas, além da possibilidade de inclusão de trabalhos e seminários no decorrer da disciplina.}

\def\PUDbibbasica{ & \href{www.biblioteca.ifce.edu.br/index.asp?codigo_sophia=24327}{Natale}, Ferdinando.  Automação industrial.  10º ed.  São Paulo, SP: Érica, 2011.  252p.  ISBN: 9788571947078.
\\ 
 & \href{www.biblioteca.ifce.edu.br/index.asp?codigo_sophia=24230}{Campos}, Mario Cesar M. Massa de.  Controles típicos de equipamentos e processos industriais.  São Paulo, SP: Edgard Blücher, 2006.  ISBN: 9788521203988. 
\\ 
 &  \href{www.biblioteca.ifce.edu.br/index.asp?codigo_sophia=23674}{Santos}, Winderson Eugênio dos.  Controladores lógicos programáveis: CLPs.  Curitiba, PR: Base Editorial, 2010.  160p.  ISBN: 9788579055737. 
}
%&  FRANCHI, Claiton Moro e CAMARGO, Valter Luís Arlindo.  Controladores Lógicos Programáveis: Sistemas Discretos.  2.  ed.  São Paulo: Érica, 2008.  ISBN 9788536501994. }

\def\PUDbibcomplementar{ & \href{www.biblioteca.ifce.edu.br/index.asp?codigo_sophia=24238}{Georgini}, Marcelo.  Automação aplicada: descrição e implementação de sistemas sequenciais com PLCs.  9º ed.  São Paulo, SP: Érica, 2007.  236p.  ISBN: 9788571947245.
%\\ 
 %&  RAMALHO, G.  L. B..  Controlador Lógico Programável.  IFCE.  2011.  56p.  ISBN 9788563953261. 
%\\ 
 %&  RAMALHO, G.  L. B..  Redes de Petri para CLP.  IFCE.  2014.  50p.  ISBN 9788563953391. 
\\ 
 &  \href{www.biblioteca.ifce.edu.br/index.asp?codigo_sophia=65382}{Capelli}, Alexandre.  Automação industrial: controle do movimento e processos contínuos.  3º ed.  São Paulo, SP: Érica, 2013.  236p.  ISBN: 9788536501178.
\\ 
 &  \href{www.biblioteca.ifce.edu.br/index.asp?codigo_sophia=63265}{Prudente}, Francesco.  Automação Industrial – PLC: Teoria e Aplicações.  2º ed.  Rio de Janeiro, RJ: LTC, 2015.  298p.  ISBN: 9788521606147.
\\
 &  \href{www.biblioteca.ifce.edu.br/index.asp?codigo_sophia=23834}{Capelli}, Alexandre.  CLP: controladores lógicos programáveis na prática.  Rio de Janeiro, RJ: Antenna, 2007.  52p.  ISBN: 9788570361370.
\\
 &  \href{www.biblioteca.ifce.edu.br/index.asp?codigo_sophia=24070}{Costa}, Cesar da.  Projetando controladores digitais com FPGA: aprenda a projetar o seu próprio Controlador Lógico Programável (CLP) utilizando FPGA de baixo custo.  São Paulo, SP: Novatec, 2006.  159p.  ISBN: 8575220888. }

\def\PUDautor{Geraldo Luis Bezerra Ramalho}

\def\PUDdata{2013-05-20}

\def\PUDrevisao{3}

\def\PUDrevisor{Antonio Barbosa de Souza Junior}

\def\PUDdatarevisao{2019-05-15}

\PUDcorpo
